% ----------------------
\subsection{Section 134 (Circa August 1835)}
% ----------------------
	\label{DC134}
	
	% -----
	\subsubsection{Historical Background}
	% -----
		\label{DC134-HB}
		
		% --
		\paragraph{JSP}
		% --
		
			``On 17 August 1835, Oliver Cowdery and Sidney Rigdon presented the Doctrine and Covenants to a general assembly of the church in Kirtland, Ohio, for approval. JS, who was visiting Michigan Territory with Frederick G. Williams, was not present at this assembly. After representatives of different priesthood offices expressed their satisfaction with the book and after `the whole congregation' unanimously `accepted and acknowledged it as the doctrine and covenants of their faith,' William W. Phelps and Cowdery presented two additional documents that the assembly `accepted and adopted, and ordered to be printed in said book': a statement on marriage and this declaration on government and law.

			``The declaration on government and law outlined the church's beliefs regarding the proper role of government in society. The most apparent motive for the declaration was the expulsion of church members from Jackson County, Missouri, in 1833. The first eleven of the declaration's twelve articles address the role and duty of government in protecting and ensuring the free exercise of religion and could be read as an indictment of local, state, and national governments for failing to fulfill their duties in protecting the church from persecution. This was not the church's first engagement in political discourse or its first reference to the constitutional rights of religious groups. In 1832, \textit{The Evening and the Morning Star} gave a detailed account of the persecutions of Quakers in colonial British North America, perhaps in an attempt to offer historical context for the growing problems faced by church members in Missouri. In March 1833, Cowdery lamented that individuals had been `persecuted for opinion's sake, or the sake of religion, when the constitution of our country allows all to worship according to the dictates of their own consciences.' Later that year, he reasserted three fundamental rights of all Americans: `the freedom of speech, the liberty of conscience, and the liberty of the press.' This August 1835 declaration contained similar language.

			``In late 1834, Missouri governor Daniel Dunklin, responding to continued pleas from church leaders for help, expressed some sympathy for the plight of church members who had been driven from Jackson County. He recognized the injustices they had suffered and wondered what laws could be amended `to guard against such acts of violence for the future.' When his statements reached Kirtland, they `revived the hopes of the church for the scattered brethren of Jackson County.' While waiting for Dunklin's sympathies to translate into action, the church initiated a series of political maneuvers. A pressing concern was to repair damage to the church's image in Missouri caused by articles in \textit{The Evening and the Morning Star} that were interpreted as supporting abolitionism and encouraging the migration of free blacks to Missouri, a pro-slavery state. To this end, in February 1835, the church founded a Democratic-leaning paper, the \textit{Northern Times}, and used it to vocally oppose abolition. Church members hoped the paper might later serve to improve relations with Democrats (including Dunklin) and pro-slavery citizens in Missouri. Although it does not denounce abolition, this August 1835 declaration also distanced the church from the slavery question, intimating a policy of noninterference. Since church leaders were making active plans in 1835 to return to Missouri in 1836, the declaration may have been seen as a formal statement of the church's position on proselytizing in slave-holding societies.

			``The declaration may have also served to quell general complaints about the church's recent political discourse. The \textit{Painesville Telegraph}, for example, criticized the church's newspaper, complaining that the church was seeking political influence. However, by publicly declaring the church's belief in constitutionally protected freedoms, the Saints' may have intended the declaration to recast their political discourse as a fundamentally American form of free speech.

			``This emphasis on American principles may have similarly helped counteract accusations that the church was anti-government. The preface to the 1835 edition of the Doctrine and Covenants states that the church had been accused of being `an enemy to all good order and uprightness' and `injurious to the peace of all governments civil and political.' However, according to the declaration, church members believed `that all men are bound to sustain and uphold the respective Governments in which they reside.' Using language and concepts from the Declaration of Independence, the Constitution, and the Bill of Rights, the declaration affirmed a dedication to law and order. It also highlighted God's hand in establishing government and emphasized the church's commitment to the separation of church and state.

			``Given that the church intended to reclaim its Missouri lands, the declaration carried special significance. To redeem Zion, the church required government support, and the declaration may have been intended to emphatically state that failure on the government's part to support the Saints in their quest was a violation of American values, constitutional rights, and the nation's founding principles.

			``The authorship of the declaration is unclear, as is the extent to which JS was involved in preparing it. Cowdery, who read the declaration at the assembly, may have been the primary author, as he editorialized on similar issues in \textit{The Evening and the Morning Star} and the \textit{Northern Times}, two periodicals for which he served as editor. Although JS's involvement in writing the declaration is unclear, he later endorsed it. The April 1836 edition of the \textit{Latter Day Saints' Messenger and Advocate} featured a letter in which JS positively cited the declaration's statement on noninterference with slaves. JS urged traveling elders to `search the Book of Covenant[s]' for `the belief of the church concerning masters and servants.' Five years later, JS signed a letter to the editor of a Chester County, Pennsylvania, newspaper in which he included the statement in its entirety, changing all of the `we believe' statements to `I believe.' Because it is uncertain how much JS contributed to the declaration's contents, the document is included as an appendix to this volume rather than as featured text.''

	% -----
	\subsubsection{Versions}
	% -----
		\label{DC134-V}
		
		See the \href{https://drive.google.com/file/d/1goAvNz8jS4R8slYfMG_db55Ty2z_vmgK/view?usp=sharing}{sections comparison} document.
	
		\begin{compactdesc}
			\item[JSP] \hyperref[EMS-1835-JS-CIR-AUG-1835-GOV]{Circa August 1835}
			\item[BC] ---
			\item[MA] 1:11 (August 1835), 163--164
			\item[DC 1835] Section 102, 252--254
			\item[DC 1844] Section 110, 440--444
			\item[DN] 2:9 (6 March 1852), 33
			\item[MS] 15:19 (7 May 1853), 300--301
		\end{compactdesc}

	% -----
	\subsubsection{Section 134:1} % *DC134-1 
	% -----
		\label{DC134-1}

		WE believe that governments were instituted of God for the benefit of man; and that he holds men accountable for their acts in relation to them, both in making laws and administering them, for the good and safety of society.

		% \begin{framed}
		% \scriptsize
			% \begin{compactdesc}
				% \item[JSP] 
				% % \item[BC] 
				% % \item[]
			% \end{compactdesc}
		% \end{framed}

	% -----
	\subsubsection{Section 134:2} % *DC134-2 
	% -----
		\label{DC134-2}

		We believe that no government can exist in peace, except such laws are framed and held inviolate as will secure to each individual the free exercise of conscience, the right and control of property, and the protection of life.

		% \begin{framed}
		% \scriptsize
			% \begin{compactdesc}
				% \item[JSP] 
				% % \item[BC] 
				% % \item[]
			% \end{compactdesc}
		% \end{framed}

	% -----
	\subsubsection{Section 134:3} % *DC134-3 
	% -----
		\label{DC134-3}

		We believe that all governments necessarily require civil officers and magistrates to enforce the laws of the same; and that such as will administer the law in equity and justice should be sought for and upheld by the voice of the people if a republic, or the will of the sovereign.

		% \begin{framed}
		% \scriptsize
			% \begin{compactdesc}
				% \item[JSP] 
				% % \item[BC] 
				% % \item[]
			% \end{compactdesc}
		% \end{framed}

	% -----
	\subsubsection{Section 134:4} % *DC134-4 
	% -----
		\label{DC134-4}

		We believe that religion is instituted of God; and that men are amenable to him, and to him only, for the exercise of it, unless their religious opinions prompt them to infringe upon the rights and liberties of others; but we do not believe that human law has a right to interfere in prescribing rules of worship to bind the consciences of men, nor dictate forms for public or private devotion; that the civil magistrate should restrain crime, but never control conscience; should punish guilt, but never suppress the freedom of the soul.

		% \begin{framed}
		% \scriptsize
			% \begin{compactdesc}
				% \item[JSP] 
				% % \item[BC] 
				% % \item[]
			% \end{compactdesc}
		% \end{framed}

	% -----
	\subsubsection{Section 134:5} % *DC134-5 
	% -----
		\label{DC134-5}

		We believe that all men are bound to sustain and uphold the respective governments in which they reside, while protected in their inherent and inalienable rights by the laws of such governments; and that sedition and rebellion are unbecoming every citizen thus protected, and should be punished accordingly; and that all governments have a right to enact such laws as in their own judgments are best calculated to secure the public interest; at the same time, however, holding sacred the freedom of conscience.

		% \begin{framed}
		% \scriptsize
			% \begin{compactdesc}
				% \item[JSP] 
				% % \item[BC] 
				% % \item[]
			% \end{compactdesc}
		% \end{framed}

	% -----
	\subsubsection{Section 134:6} % *DC134-6 
	% -----
		\label{DC134-6}

		We believe that every man should be honored in his station, rulers and magistrates as such, being placed for the protection of the innocent and the punishment of the guilty; and that to the laws all men show respect and deference, as without them peace and harmony would be supplanted by anarchy and terror; human laws being instituted for the express purpose of regulating our interests as individuals and nations, between man and man; and divine laws given of heaven, prescribing rules on spiritual concerns, for faith and worship, both to be answered by man to his Maker.

		% \begin{framed}
		% \scriptsize
			% \begin{compactdesc}
				% \item[JSP] 
				% % \item[BC] 
				% % \item[]
			% \end{compactdesc}
		% \end{framed}

	% -----
	\subsubsection{Section 134:7} % *DC134-7 
	% -----
		\label{DC134-7}

		We believe that rulers, states, and governments have a right, and are bound to enact laws for the protection of all citizens in the free exercise of their religious belief; but we do not believe that they have a right in justice to deprive citizens of this privilege, or proscribe them in their opinions, so long as a regard and reverence are shown to the laws and such religious opinions do not justify sedition nor conspiracy.

		% \begin{framed}
		% \scriptsize
			% \begin{compactdesc}
				% \item[JSP] 
				% % \item[BC] 
				% % \item[]
			% \end{compactdesc}
		% \end{framed}

	% -----
	\subsubsection{Section 134:8} % *DC134-8 
	% -----
		\label{DC134-8}

		We believe that the commission of crime should be punished according to the nature of the offense; that murder, treason, robbery, theft, and the breach of the general peace, in all respects, should be punished according to their criminality and their tendency to evil among men, by the laws of that government in which the offense is committed; and for the public peace and tranquility all men should step forward and use their ability in bringing offenders against good laws to punishment.

		% \begin{framed}
		% \scriptsize
			% \begin{compactdesc}
				% \item[JSP] 
				% % \item[BC] 
				% % \item[]
			% \end{compactdesc}
		% \end{framed}

	% -----
	\subsubsection{Section 134:9} % *DC134-9 
	% -----
		\label{DC134-9}

		We do not believe it just to mingle religious influence with civil government, whereby one religious society is fostered and another proscribed in its spiritual privileges, and the individual rights of its members, as citizens, denied.

		% \begin{framed}
		% \scriptsize
			% \begin{compactdesc}
				% \item[JSP] 
				% % \item[BC] 
				% % \item[]
			% \end{compactdesc}
		% \end{framed}

	% -----
	\subsubsection{Section 134:10} % *DC134-10 
	% -----
		\label{DC134-10}

		We believe that all religious societies have a right to deal with their members for disorderly conduct, according to the rules and regulations of such societies; provided that such dealings be for fellowship and good standing; but we do not believe that any religious society has authority to try men on the right of property or life, to take from them this world's goods, or to put them in jeopardy of either life or limb, or to inflict any physical punishment upon them.  They can only excommunicate them from their society, and withdraw from them their fellowship.

		% \begin{framed}
		% \scriptsize
			% \begin{compactdesc}
				% \item[JSP] 
				% % \item[BC] 
				% % \item[]
			% \end{compactdesc}
		% \end{framed}

	% -----
	\subsubsection{Section 134:11} % *DC134-11 
	% -----
		\label{DC134-11}

		We believe that men should appeal to the civil law for redress of all wrongs and grievances, where personal abuse is inflicted or the right of property or character infringed, where such laws exist as will protect the same; but we believe that all men are justified in defending themselves, their friends, and property, and the government, from the unlawful assaults and encroachments of all persons in times of exigency, where immediate appeal cannot be made to the laws, and relief afforded.

		% \begin{framed}
		% \scriptsize
			% \begin{compactdesc}
				% \item[JSP] 
				% % \item[BC] 
				% % \item[]
			% \end{compactdesc}
		% \end{framed}

	% -----
	\subsubsection{Section 134:12} % *DC134-12 
	% -----
		\label{DC134-12}

		We believe it just to preach the gospel to the nations of the earth, and warn the righteous to save themselves from the corruption of the world; but we do not believe it right to interfere with bond-servants, neither preach the gospel to, nor baptize them contrary to the will and wish of their masters, nor to meddle with or influence them in the least to cause them to be dissatisfied with their situations in this life, thereby jeopardizing the lives of men; such interference we believe to be unlawful and unjust, and dangerous to the peace of every government allowing human beings to be held in servitude.

		% \begin{framed}
		% \scriptsize
			% \begin{compactdesc}
				% \item[JSP] 
				% % \item[BC] 
				% % \item[]
			% \end{compactdesc}
		% \end{framed}